\section{Scheduling Passes}
\label{sec:method}

% VALUE: the passes and their legality checkers are the implementation: the composer reorders and renames across the operators of the open prefill, the stream pass reorders within one operator's stream, the loop pass renames inside decode attention's loop, and each output is checked for legality from the instructions and run bit-identically; the event model makes each choice before anything runs
Our passes remove false dependences of the classes of Section~\ref{subsec:grid}, and legality checkers accept each rewrite before it runs.
In decode's loop the class and the buffers come from the analysis of Section~\ref{subsec:reading}; in the CP's stream they were identified by reading the designs, and the analysis supplies the graph against which each rewrite is checked.
For the open prefill, a composer takes an existing design's per-operator instruction streams, its map of operators to columns, the dependency graph of its host program, and the probes, and emits a sequence of composed commands.
The composed commands must do the same work as the existing ones, which the composer guarantees by keeping every instruction and checking legality from the instructions themselves (Section~\ref{subsec:compose}).
No group may wait at the CP for work it does not consume, which is an order dependence in the stream row of Table~\ref{tab:grid}, so the composer reorders issue across operators and a stream pass reorders issue within one operator's stream (Section~\ref{subsec:order}).
No use of a buffer may wait for another use it does not depend on, which is a name dependence in the same table, so the composer renames buffers into rings and our loop pass renames decode attention's buffers by the same rule (Section~\ref{subsec:rings}).
The partition of the array must fit the regime, so the composer chooses the split of columns for each regime (Section~\ref{subsec:split}).
The event model of Section~\ref{subsec:pricing} prices every candidate, so each choice is made without running the candidate.
Removing a dependence is not always worth its cost, and Section~\ref{subsec:removal} gives both the rule that decides and the regime where the composed order loses time.
% SOURCE: amd-npu a0ee2edf44e0 (tag composer-v1.1), prefill-overlap/txn.py, merged.py, component_model.py; Python, on the engine's hardware context; the stream pass aie-cosched-schedule and the loop pass aie-cosched-rename in mlir-aie (aie-cosched-pass-0 f1725dae1f1).

\subsection{Composing Instruction Streams}
\label{subsec:compose}

% VALUE: a composed command is an interleaving of operators' segments that keeps each operator's own order and every instruction once, so only the timing of issue and waits changes
The composer splits each operator's instruction stream, already reduced to its own columns, into the segments of Section~\ref{subsec:formulation}.
A composed command concatenates segments of several operators in a chosen order, keeping each operator's segments in their own order and every instruction exactly once.
Under these rules, an interleaving changes only when each operator's work is issued and awaited, never what the work is.
% SOURCE: txn.segments and txn.compose (amd-npu a0ee2ed prefill-overlap/txn.py); compose asserts each stream's segment order, every op kept once, and equal headers.

% VALUE: two operators may interleave only if they cannot disturb each other's device state, which the composer checks statically from the instructions themselves; operators on one group stay in sequence
Two operators may interleave in one command only if they cannot disturb each other's state on the device, which the composer checks from the instructions themselves.
Every register write of an operator must address a column of its own group and touch only shim BDs and DMA queues, the operator's own cores' locks, and runtime parameters in its own cores' memory.
Interleaved operators must patch disjoint argument slots or disjoint byte ranges of one slot's buffer, and no operator may write a buffer that another operator in the same command reads or writes.
Operators on the same group, such as several GEMMs, are composed in sequence rather than interleaved, so their writes to that group's registers keep the order of separate commands.
% SOURCE: each check turns red on a defeated input (swapped groups, an unstripped stream, a synchronization before an issue op, a patch onto another group's argument, a duplicated column, a reordered or dropped segment) (lab-prefill-overlap (3), (4)).

\subsection{Reordering Issue}
\label{subsec:order}

% VALUE: the composer list-schedules the dependency graph chunk by chunk rather than stage by stage, so each attention operator meets independent GEMMs of earlier chunks, and it sizes each command so its GEMM work matches the attention work
The existing host program runs each layer stage by stage, so every output projection waits behind the layer's entire attention, and no GEMM is ready to run beside attention.
The composer instead list-schedules the dependency graph chunk by chunk.
At each step, it takes the ready attention operator of the lowest layer and chunk and adds ready GEMMs, in the same order, until their estimated work reaches that of the attention operator.
The GEMMs then run in sequence on columns 1--6 while attention runs on columns 0 and 7.
% VALUE: the price of running the two groups at once: the columns share the DRAM read cap, so a group whose partner is bandwidth-bound loses a measured share of its partner's time, and the ceiling on a pair is below the ideal; where a group is already near the cap there is nothing to win
Running the two groups at once is not free, because the columns share one DRAM read rate.
A group on the critical path loses 8.5\% of its partner's duration when it is bandwidth-bound and almost nothing when it is not, so two operators of equal duration overlap to within $1.84\times$ of the longer rather than the ideal $2.00\times$ (Section~\ref{subsec:removal}).
Where a design already reads near the cap the reorder therefore has little to recover, which Section~\ref{subsec:bound-eval} reports for the largest model we evaluate.
% SOURCE: the penalty by regime and the 1.84x ceiling are Section 5.1's measurement; per command on Llama-3.2-1B at 4,096 tokens, commands carrying GEMMs read 90-93% of the cap and commands carrying attention about 15% (vault decisions log, "Regime analysis" steps 1 and 3); the 8B composed prefill averages 93-96% of the cap (lab-metric-fields (1)).
% DECISION: tradeoffs are stated beside the decision that incurs them rather than collected in a section (author, 2026-10-21).
For an append to a cached context, the composer plans only the new chunks, and their attention estimates use each chunk's absolute index.
% SOURCE: merged.plan (a0ee2ed): qkv(L, c) after down(L-1, c); attn(L, c) after qkv(L, 0..c); o after attn; gate and up after o; down after gate and up; an optional cap on the GEMMs per attention command.
% SOURCE: append = Merged.prefill(ids, start=L) on the same plan rules; attention estimates linear in the absolute chunk index (vault lab-agent-step (1)).

% VALUE: within a command, the CP opens every group's work first and then places each synchronization at its predicted completion, followed by the issue it gates, so no synchronization holds back another group's work; the order is fixed before dispatch from the estimates
Within a command, the order of synchronizations decides how long each group waits for its next issue (Figure~\ref{fig:order}).
In turn, the CP would issue the first GEMM, then walk the attention stream to its end, and only then issue the second GEMM, so the GEMM group idles behind attention's synchronizations.
The composer instead issues the opening segments of both attention columns and of the first GEMM, and then places every synchronization in order of its \emph{predicted completion}, each followed by the issue it gates.
A GEMM's synchronizations are predicted to complete after the sum of the estimates of the GEMMs before it, and attention's mid-stream synchronizations at fixed fractions of its estimate.
Attention's two columns are treated as two groups, so their segments alternate, which removes the false dependence inside the attention operator (Section~\ref{subsec:losses}).
The order is static: the estimates fix it when the command is built, and the device executes it without host involvement.
When issue costs are negligible, the estimates are exact, and the host work is hidden, the CP never blocks on a later-completing operator while an earlier one's successor waits, so the order matches the list schedule of the groups without a CP.
% SOURCE: merged.cp_order (a0ee2ed); mid-stream fractions (0.25, 0.5) of attention's estimate, the lowest of five placements tried; the in-turn order takes 1.26x as long on the same 38 commands (lab-prefill-overlap (16)).
% DECISION: the exactness statement in the last sentence is argued, not proven; it stays a remark until a proof is written (vault ch3-formulation, Lower bound and heuristic).

% VALUE: within one operator's stream, the stream pass's full mode interleaves the columns' segments round-robin by index, the rule that produced the hand-alternated streams, so it emits the alternation byte-identically on the population's attention streams, and the legality checker accepts each rewrite by checking every inverted pair against the decoded configuration
Within one operator's stream, a stream pass on the lifted runtime sequence reorders issue as well.
In its full mode, it interleaves the columns' segments round-robin by index, which is how the alternated attention streams were produced by hand, so it emits that alternation rather than discovering it.
On all 42 attention streams of the population, the pass's output is byte-identical to the hand-alternated stream.
Reproducing a hand edit is the point rather than the limit: a hand edit reaches the streams someone reads, whereas a pass that must accept or refuse every stream has to carry the condition under which the rewrite is legal, which is what lets it reach the streams nobody has read.
A legality checker accepts each rewrite before it runs, by checking every pair of instructions that the rewrite inverts against the routes, BDs, and buffer ranges decoded from the configuration.
% SOURCE: aie-cosched-schedule (mlir-aie 573e0140174, identical in f1725dae1f1), proposal format /2 frozen at aie-cosched 8dee63f; 42 / 42 decoded order and bytes equal to the hand versions (vault lab-cosched-passes (35), (40)); the pairwise inversion rule of legality.py (lab-cosched-passes (14)).
% DECISION: the pass emits the alternation and the checker accepts it; it is never said to discover it, and full mode is never called a predicted-completion rule (Ruling D3, vault decisions log).

\subsection{Renaming Buffers}
\label{subsec:rings}

% VALUE: with commands in flight and chunks interleaved, two uses of one buffer would wait on each other although neither consumes the other's data; the composer gives each overlapping use its own slot, with ring depths derived from the plan's live intervals and a check that no slot is reused before its reader is done
Once operators of several chunks are in flight, a GEMM's input and output buffers cannot be reused per stage as in the existing host program, because two uses of one buffer would then wait on each other.
For each buffer role, the composer derives a ring of slots from the plan: each use of a buffer lives from its producer's write to its last reader's completion, and interval colouring assigns slots so that no two overlapping uses share one.
A separate check asserts, over the whole plan, that no slot is reused before its previous reader has completed.
The depths follow from the plan and are not set by hand; for example, every QKV input of a layer is written before the layer's first QKV command completes, so that ring needs one slot per chunk.
% VALUE: the price of the rings: slots are memory, and the composer buys overlap with it; the arenas replace the engine's own weight buffers rather than adding to them, so the cost is the schedule's slots and not a second copy of the weights
Slots are memory, so this is where the composed schedule pays for its overlap.
The rings and the plan's own buffers take 53--179\,MB across the configurations we build, whereas the arenas that hold the weights replace the engine's own weight buffers rather than adding to them, so the weights are not copied twice.
Peak live device memory for a 4{,}096-token append to a 16{,}384-token context is then 3.5\,GB on Llama-3.2-1B, 7.7\,GB on Qwen3-4B, and 9.7\,GB on Llama-3.1-8B; the production engine holds 2.2 and 7.9\,GB for the first two on the same weights, so the schedule's slots and not a second copy of the weights are what we add.
% SOURCE: composer schedule buffers (Layout plus lazy) 53.4-178.8 MB over the 21 configurations (min 1B S4 16k + 4k, max 8B S2 scratch 2k; amd-npu 71901c50 metrics/retro-2026-10-21/footprint-summary.txt, re-derived in the 2026-10-22 foundation audit); weight arenas 645 MB (1B), 2,340 (Qwen3-4B), 4,451 (8B), replacing rather than duplicating the engine's weight buffers, with peak minus live equal to the arena size within 2% at load; peak live 2,206 / 5,923 / 7,399 MB at 4,096 tokens from scratch and 3,471 / 7,725 / 9,721 MB at 16,384 + 4,096; FastFlowLM's server at drm-total-memory 2,110,100 KiB (1B) and 7,668,316 KiB (Qwen3-4B), i.e. 2.16 and 7.85 GB (mapped 2,158,084,096 and 7,847,751,680 B; the earlier 2.11 / 7.67 read KiB as kB) (vault lab-metric-fields (1); lab-flm-practice (3)).
% DECISION: the tradeoffs are stated beside the decision that incurs them rather than collected in a section (author, 2026-10-21).
% SOURCE: merged.colour, assert_no_hazard, and Layout (a0ee2ed); the hazard check turns red on an overlapping slot; ring depths at 4,096 tokens: shared 2048-wide input 32, 8192-wide input 4, outputs 4-18 (lab-prefill-overlap (9)).

% VALUE: decode attention holds the same false dependence inside one operator: one buffer each for a block's keys, values, and scores makes each block wait for the previous one; renaming alone gains nothing while the scores are acquired once per pair of blocks, as in a hand edit unrolled by two, because the two cores still wait for each other; our loop pass renames all three into two-slot rings acquired once per block on both ends, keeps the carried state, and matches a hand edit's gain with unchanged kernels and bit-identical outputs
Decode attention holds the same false dependence inside one operator, where a single KV buffer makes each block's transfer wait for the compute on the previous block (Section~\ref{subsec:losses}).
Only the running maximum and the accumulator must carry across blocks, whereas each block's keys, values, and scores are written once and read once, so a second buffer for each lets transfer and compute alternate between them.
Renaming alone is not enough, however, when the scores are acquired once per pair of blocks, as in a hand edit that unrolls the loop by two.
With the keys, values, and scores renamed but the scores acquired per pair, the time per token stays within 2\% of the existing design's, since the core that computes a block's scores and the core that folds them into the output still wait for each other.
Our loop pass therefore renames each such buffer into a ring of two slots that each core acquires and releases once per block, so the first core computes block $b+1$ while the second folds block $b$.
The compiler's double-buffering transformation declines the loop because of the carried state, whereas our pass classes every buffer of the lowered loop by its conflicts across iterations and leaves the carried state in place; the kernels and the arithmetic order are unchanged, and the decoded tokens and logits stay bit-identical.
On both models, the pass's output comes within 1\% of a hand edit that passes the scores between the cores through a channel instead, at every evaluated length.
% SOURCE: renamed K, V, and score in the hand edit unrolled by two, scores acquired per pair (arm (b)): 0.997-0.984x per token at 4k-32.5k, slope 1.181 vs 1.153 us (Llama-3.2-1B; vault lab-decode-long-context (13), amd-npu aab2aae); the same edit on Llama-3.1-8B 0.990x at both cells (lab-decode-long-context, decode nulls Phase 2). Loop pass (aie-cosched-rename, mlir-aie 41b4947f09e; classes and rings in vault lab-cosched-passes (33)): against the existing design 1.269x at 16,384 and 1.391x at 32,512 (1B), 1.160x at 16,384 and 1.205x at 28,800 (Qwen3-4B); pass / hand edit 0.992-1.004 at every point; slope 0.683 us (1B) against a transfer-only 0.677 us; tokens and every step's logits bit-identical across the existing design, the hand edit, and the pass, 1B P = 0 to 32,517 and Qwen3-4B P = 2,048 to 28,800 (vault lab-decode-long-context (22), amd-npu a6bc94b and 455740b).
% DECISION: the rename is our loop pass's output on the lowered design, not a hand edit; the hand edit appears only as the reference the pass reproduces.

% MOVED from the former limits paragraph: the scope of what these passes touch
Our passes reorder and rename within one configuration of the array and leave every reconfiguration point where the stack put it, so the gains we report hold with those points unchanged and say nothing about moving them.

\subsection{Partitioning per Regime}
\label{subsec:split}

% VALUE: the split of columns between attention and GEMMs is a build-time parameter of the existing design, fixed by hand; the legal splits follow from the heads, each split is a separate configuration built in seconds, and the composer needs only more argument slots to serve it
The existing design fixes the split of columns between attention and the GEMMs by hand: attention groups are listed per column, and the GEMMs take a contiguous block of the remaining columns.
Since each attention group serves an equal share of the KV heads, a model with 8 KV heads admits attention on 1, 2, or 4 columns, with the GEMMs on the other 7, 6, or 4.
We build each split as its own configuration from the unchanged design source, which takes seconds, and each split's outputs are bit-identical to those of the default split.
The composer serves every split with three argument slots per attention group instead of a fixed six, and it checks every GEMM's encoded iteration count and stride against the shim BD's field limits before it uses a split.
% SOURCE: mlir-air a3e26261 fused_prefill/device.py: Attn(col=...) per group in Config.attn, GCOL0, GNC = 1, 6 as module constants; spec() splits kv_heads evenly, so 1, 2, or 4 attention columns with 8 KV heads (vault lab-split-sweep (1)).
% SOURCE: out-of-tree split modules, builds of 13-14 s; logits bit-identical to the default split in every configuration on both models (vault lab-split-sweep (2), (4), (5)).
% SOURCE: BD iteration wrap 6 bits (at most 64 rounds), stride 20 bits (words); the four-column split's LM-head GEMM runs 48 rounds; decoded from aie-rt xaie2pgbl_params.h fields with a round-trip check (vault lab-split-sweep (5)).

% VALUE: the composer picks the split per regime with the event model from one probe per split; more attention columns shorten attention, whose time grows with the context, but lengthen every GEMM, so the choice depends on the prompt and the cached context
For each split, the composer probes every operator alone once and predicts the TTFT of each regime with the event model (Section~\ref{subsec:pricing}), and it picks the split with the lower prediction.
More attention columns shorten attention, whose time grows with the context, but lengthen every GEMM, so the choice depends on the prompt length and on the cached context.
Which group binds the lower bound decides the direction: the GEMMs bind for prompts computed from scratch, so the faster split gives them more columns, whereas attention binds for appends to a long cached context, so the faster split gives attention more (Section~\ref{subsec:bound-eval}).
% SOURCE: per-split probes at n = 2,048, prediction with component_model.simulate (amd-npu 0f61651 prefill-overlap/g1_split_model.py, vault lab-split-sweep (6)-(7)).
% SOURCE: the four-column split's GEMMs are 27% (1B) and 35% (Qwen3-4B) slower than on six columns, and its attention slope per chunk 3.0-3.1x lower than the default split's existing stream (vault lab-split-sweep (4)-(5)).

% VALUE: the chosen splits stay resident, so a conversation switches between them at a turn boundary by a suspend and resume; its cached context, laid out by the attention groups of the split that computed it, moves with it
A deployment keeps each chosen split resident in its own hardware context, with its own copy of the weights, and switches between them at a turn boundary by suspending one context and resuming the other.
The KV cache of a cached context is laid out by the attention groups of the split that computed it, so a conversation that changes split also moves its cache to the new layout.
% VALUE: the price of making the partition a per-regime choice: each split is its own configuration, so a deployment either holds two of them with two copies of the weights or pays a switch and a cache relayout; the weights are the larger of the two costs on a large model
Making the partition a choice therefore costs either memory or a switch.
Each split is its own configuration, so holding both resident holds the weights twice, which on the largest model we evaluate is 4.5\,GB of weight arenas for each; releasing one instead costs a warm switch, 7--8\ms{} on the two smaller models, and the cache relayout above.
Neither cost applies within one regime, where a single split serves every turn.
% SOURCE: weight arenas 645 MB (Llama-3.2-1B), 2,340 (Qwen3-4B), 4,451 (Llama-3.1-8B), one copy per resident split (vault lab-metric-fields (1)); warm switch 7.1 ms (1B) and 7.7 ms (Qwen3-4B), suspend and resume only, without moving a cached context (lab-split-sweep (5)).
Moving the cached context between splits is the cost we do not measure, and it is the one that decides between the two paths.
Its floor, reading and writing the per-block records once at the part's peak memory bandwidth, is 0.6--4.1\,ms on Llama-3.2-1B and 2.7--18.6\,ms on Qwen3-4B over cached contexts of 4{,}096 to 28{,}672 tokens; the closest host operation we have measured, the scatter that seeds a decoder, puts it nearer 26--181\,ms and 116--814\,ms.
On the measured figure the move costs between 4 and 106 times the warm switch, so a deployment that switches regimes often holds both splits resident and pays the memory instead; on the floor alone the smallest model's move would be cheaper than the switch, so we state both and claim neither as the price.
% SOURCE: computed host-only by npu-gap and recorded here: records are layers x kv-heads x ceil(L/32) blocks at 4,608 B (Llama-3.2-1B) and 9,216 B (Qwen3-4B) per block per head, giving 75.5 / 302.0 / 528.5 MB and 339.7 / 1,359.0 / 2,378.2 MB at 4k / 16k / 28k. Floor = 2R at the part's 256 GB/s peak (lscpu, a specification figure, not measured). The nearer figure applies the restack fit's seed term, 342.16 ms per GB written (amd-npu a61417d0 fit-2026-10-21.json), which is an order of magnitude and not a price: it includes a cast a record copy does not do, excludes the per-block sync, and was measured in the clean host state, the same work having fallen into the huge-page excursion of Section 8.1 at 6-14x.
% DECISION: npu-gap's bound contradicted my expectation that the move is an order of magnitude above the switch everywhere; it is not, on the floor at the 1B. Both numbers are stated and neither is called the price.
% SOURCE: switch = suspend + resume + first prefill in one process, both splits' weights resident (vault lab-split-sweep (5)).
% DECISION: the relayout and its decision rule enter once measured (vault paper-cosched-plan, regime switching); until then this sentence stays a placeholder.

% VALUE: what a platform must offer for the method to apply: an in-order issuer whose issue returns at once and whose waits block, disjoint groups that share only DRAM, per-operator streams, and buffers named by slot and offset
The composer relies on four properties of the platform rather than on MLIR-AIR: a single in-order issuer whose issue instructions return at once and whose synchronizations block, column groups that share only DRAM bandwidth, a design whose operators expose their own instruction streams, and buffers that commands name by an argument slot and a byte offset.
% SOURCE: Sections 2.1 and 3.1; the five q4nx models of the evaluation all run on AIR's shared fused prefill and fused decode (mlir-air a3e26261 programming_examples/llms/*_q4nx).

\subsection{Implementation}
\label{subsec:impl}

% VALUE: the composer is a small Python layer that replaces the host program's dispatch loop on the engine's own hardware context; it changes no kernel, compiler, driver, or firmware
The composer replaces the dispatch loop of the design's host program and submit to the engine's own hardware context, with no change to a kernel, the compiler, the driver, or the firmware.
% SOURCE: amd-npu a0ee2edf44e0 (tag composer-v1.1), prefill-overlap/txn.py (111 lines), merged.py (616), component_model.py (180).

\textbf{Argument Arenas.}
% VALUE: the existing design binds one buffer per argument slot, so a command holds only one GEMM; we retarget only a patch's byte offset, never its slot, so any number of GEMMs share the GEMM slots through arenas, under an offset bound the composer respects
The existing design gives each operator kind its own argument slots, so a command that carries attention can hold only one GEMM.
Since retargeting a patch to another slot faults in the IOMMU whereas changing its byte offset is safe, we allocate the GEMM inputs, weights, and outputs as sub-buffers of three parent buffers, which we call \emph{arenas}, and rewrite only each GEMM's patch offsets, so any number of GEMMs share three slots.
The offset field bounds an arena to $2^{31}$ bytes, so the weights are split into arenas of whole layers, and the GEMMs of one command share a weight arena.
% SOURCE: nine slots, nargs = 3 + 3 x len(cfg.attn) (mlir-air a3e26261 engine.py:107); address patch = opcode 0x81, argument index at word 8, byte offset at word 10.
% SOURCE: slot moves complete with wrong output and AMD-Vi IO_PAGE_FAULT lines; offset moves inside a parent buffer are bit-identical; two GEMMs at distinct offsets in shared arenas are bit-identical (amd-npu 61c64f9 prefill-overlap/retarget/, lab-prefill-overlap (8)).
% SOURCE: merged.Weights (a0ee2ed); Llama-3.2-1B uses one weight arena, Qwen3-4B two (2145 + 195 MB) (lab-prefill-overlap (24), (26)).

\textbf{Host Boundaries.}
% VALUE: host work between commands forces boundaries; keeping D commands in flight hides the host work and the boundary cost but delays which operators are ready, so the depth and the packing of GEMMs without attention are chosen per prompt length with the event model
Since the host computes normalizations, rotary embeddings, and residuals between commands, a dependence that crosses the host forces a command boundary.
The composer keeps $D$ commands in flight and processes a command's outputs while later commands run, so a command can carry only operators whose inputs come from commands at least $D$ positions earlier.
A step without a ready attention operator packs a fixed number of ready GEMMs into one command, following standard practice for amortizing dispatch cost (Section~\ref{subsec:semantics}).
The event model chooses the depth and the packing for each prompt length, and we report their effect apart from that of the essential decisions (Section~\ref{subsec:formulation}).
% VALUE: the price of depth: it hides the host's work but delays which operators are ready, and what it is worth falls as the prompt grows, so the right depth is a property of the workload rather than of the design
Depth trades the host's exposed time against how early an operator becomes ready, and what it buys falls as the prompt grows: at a depth of two the host is exposed for 9.3\% of TTFT at 2{,}048 tokens and for almost none at 4{,}096.
The existing host program instead keeps every command of a stage in flight, 32--64 of them, and so pays no exposed host time and gains no choice of order.
% SOURCE: host-exposure share under the composer's policy 0.5% of TTFT at 2k and about 0 at 4k, against 9.3% without queue depth at 2k; the shipped program keeps 32-64 commands in flight against the composer's 2 (vault lab-prefill-overlap (16) and its pre-registration; lab-split-sweep (18)).
% SOURCE: merged.plan with lag = D; Layout claims an output slot when its command is queued, D - 1.6 commands before it runs (a0ee2ed).

\textbf{Decode Loop Renaming.}
% VALUE: what our loop pass changes in the lowered decode design, in local memory and in the configuration, what it skips and why, and that it keeps one kernel call per block and leaves every build output unchanged when switched off
The pass gives each renamed buffer a twin on the same tile, raises its free-slot lock from one slot to two, extends its buffer-descriptor chain through the twin, and selects the slot in each iteration through a one-word rotation counter in local memory that the configuration initializes.
It keeps one kernel call per block, so the loop runs exactly the blocks that the context needs, and it lists with a reason, rather than renames, any buffer whose synchronization is not one lock pair or whose twin would exceed the tile's 16 buffer descriptors or 64\,KiB of local memory.
The pass runs as an opt-in step inside the existing compiler, and with it switched off, every instruction stream and configuration that the compiler emits is byte-identical to the existing build.
% SOURCE: aie-cosched-rename, mlir-aie 41b4947f09e: 12 buffers renamed on the 1B census build (K, V, score per attention pair), 0 over budget, 16 / 16 rotation words written by the PDI's initial-value writes (vault lab-cosched-passes (33)); hook-off identity PASS on the 1B and Qwen3-4B, differences confined to xclbin / PDI UUIDs and timestamps, as between two stock builds (vault lab-decode-long-context (20)-(22), control 2).
